% ============================================================================
% v5 front matter -- reframed. Drop-in replacement for main.tex lines 34-113
% (title, abstract, keywords, Section I). Everything else migrates per
% paper/V5_PLAN.md.
%
% Thesis: evaluation practice for forecast-driven energy management cannot
% detect total failure. Continual learning is the system under test, not the
% claim.
%
% The ledger linearity, the concavity derivation, the SAA stochastic program
% and the RL comparison are deliberately ABSENT: they belong to paper 2. This
% paper diagnoses the truncation defect and leaves it open.
% ============================================================================

\title{Reported versus Delivered Savings: A Deployment Audit of
Forecast-Driven EV Charging Control}

\author{
\IEEEauthorblockN{
Abdelhaq Dabaghia\IEEEauthorrefmark{1}\IEEEauthorrefmark{2}
and Pedro Rodriguez\IEEEauthorrefmark{1}\IEEEauthorrefmark{2}
}
\IEEEauthorblockA{
\IEEEauthorrefmark{1}
Luxembourg Institute of Science and Technology (LIST)\\
L-4362 Esch-sur-Alzette, Luxembourg\\
Email: abdelhaq.dabaghia@list.lu, pedro.rodriguez@list.lu
}
\IEEEauthorblockA{
\IEEEauthorrefmark{2}
Department of Engineering, University of Luxembourg\\
L-4365 Esch-sur-Alzette, Luxembourg
}
}

\maketitle

\begin{abstract}
Forecast-driven energy management systems are routinely reported through
savings that are computed rather than measured. We deploy a continual-learning
forecasting and model-predictive-control system at a commercial site in
Luxembourg, instrument it to settle its own economics against metered demand,
and report what the instrumentation found. Over a fortnight of operation the
controller computed schedules, wrote key-performance indicators claiming
savings, and dispatched nothing: every publish to the chargers failed, 400
consecutive times, with no successes. The same period ran on a constant
electricity price, because the market-data credential had been rejected for ten
weeks and the price source degraded silently to a flat default---under which an
energy-conserving shift cannot change the bill by construction. Seven
independent components were found answering plausibly instead of refusing.
Re-scoring the unchanged production weights under the feature convention
actually served raises retained-task error from $0.548$ to $1.254$~nRMSE and
drops correlation with ground truth from $0.958$ to $0.458$; a causal refit
recovers most of it. That correction also withdraws one of our own published
contributions: the proximal regulariser, which beat fine-tuning under the leaky
pipeline, is indistinguishable from it under causal features ($p=0.6953$). What
survives is stronger---bounded rehearsal beats fine-tuning by $0.346$~nRMSE
unanimously over ten seeds ($d_z=-17.57$) and is the only mechanism to beat
naive persistence---yet that accuracy advantage does not reach the decisions it
exists to support: forecasts $34\%$ better on all $35$ day--seed pairs leave
realised controller cost statistically unchanged ($p=0.645$). We contribute a
settlement methodology separating planned, counterfactual and settled savings,
a measured account of silent no-op degradation, and an open guard suite mapping
each defect class to an executable test.
\end{abstract}

\begin{IEEEkeywords}
Evaluation validity, deployment audit, settlement, train/serve skew, feature
leakage, continual learning, model predictive control, electric vehicle
charging, commercial energy site
\end{IEEEkeywords}

% ============================================================================
\section{Introduction}
\label{sec:intro}

Commercial energy sites increasingly couple machine-learned forecasts to
optimisation-based controllers: a model predicts photovoltaic generation and
electric-vehicle charging demand, and a scheduler turns those predictions into
setpoints that move real power. The literature on such systems is large and
overwhelmingly positive, and it is reported almost entirely through a single
number---the saving the controller expects to make.

That number is a \emph{plan}. It is computed by subtracting the cost of an
optimised schedule from the cost of an unmanaged one, with both terms evaluated
on the same forecast. Nothing in the arithmetic touches a meter. A system whose
setpoints never reach the hardware, whose price feed has silently frozen, or
whose forecaster is being served features it was not trained on, produces
exactly the same number as one working perfectly.

We know this because we built such a system, deployed it at a commercial site
in Luxembourg, and then instrumented it to settle its own economics against
metered demand. The instrumentation is what this paper reports.

What it found was not a degradation. For a fortnight the controller solved its
optimisation on schedule, wrote key-performance indicators claiming savings,
and dispatched nothing at all: the publish to the chargers failed on every
cycle, 400 consecutive times, with zero successes. One of the two metering
points was never addressed by the control path in the first place. Over the
same period the scheduler was optimising against a constant price, because the
market-data credential had been rejected for ten weeks and the price client
degraded to a flat default---and at a flat price an energy-conserving load
shift cannot change the bill, so the true saving was exactly zero every cycle.
None of this appeared in any log as an error, any dashboard as a gap, or any
alarm as a page.

The pattern repeated. A promotion gate that compared candidates only to their
predecessor, never to a naive baseline, and would therefore promote a model
that persistence beats by a factor of $2.3$. An ingest logging \texttt{0 rows
upserted} as success, so that a dead feed froze telemetry for two days
invisibly. A lookback anchored to the present and filled with the training
median wherever data had not arrived, so that roughly a fifth of the window,
and specifically its most recent fifth, was fabricated. In each case the
component answered plausibly instead of refusing, and in most of them the
failure was found because a person thought a display looked wrong.

The same audit reached the model itself. The endogenous rolling means used for
training were computed over windows containing the prediction target, while the
deployed service computed them causally; re-scoring the unchanged production
weights under the convention actually served raises retained-task error from
$0.548$ to $1.254$~nRMSE. This is not only a deployment defect. It reverses a
published conclusion of our own: a proximal regulariser that beat naive
fine-tuning under the leaky pipeline is statistically indistinguishable from it
once features are made causal. We report that withdrawal here, with the
evidence that forces it.

The contributions of this paper are therefore about measurement rather than
method. It proposes no new learning algorithm, and says so plainly.

\begin{enumerate}

\item A \textbf{settlement methodology for forecast-driven site control},
separating \emph{planned} savings (forecast against forecast), \emph{counter-
factual} savings (a plan settled against metered demand it never actually
drove) and \emph{settled} savings (a dispatched plan settled against the
demand that arrived), together with the actuation-record boundary that
determines when the control database may be trusted about what was commanded.

\item A \textbf{measured account of silent no-op degradation}: seven
independent components at one site that answered plausibly instead of
refusing, tabulated with the mechanism, the duration before detection, and---
the column that matters---how each was eventually found.

\item \textbf{Train/serve feature skew, quantified by controlled ablation on a
live system}: the same weights on the same holdout, re-scored under the served
convention, together with the causal refit that recovers most of the loss.

\item A \textbf{falsification of one of our own contributions}, with the
statistics that force it, and the result that survives correction: bounded
rehearsal dominates both naive fine-tuning and the regulariser, unanimously
over ten seeds, and is the only mechanism to beat naive persistence.

\item \textbf{Evidence that forecast accuracy does not predict decision
value}, together with a bound on what any online arbiter could recover from
choosing between adaptation mechanisms.

\end{enumerate}

We also release the instrumentation: a guard suite that maps each defect class
identified here to an executable test, so that the failures described are
detectable by others rather than merely narrated.

Section~\ref{sec:related} positions this against work on deployment reporting
and evaluation validity. Section~\ref{sec:arch} describes the site and the
deployed architecture, and Section~\ref{sec:cl} the continual-learning system
under test. Section~\ref{sec:mpc} gives the control formulation and
Section~\ref{sec:settle} the settlement methodology.
Section~\ref{sec:defects} presents the defects and what each cost.
Section~\ref{sec:setup} and Section~\ref{sec:results} give the experimental
protocol and results, Section~\ref{sec:discussion} the limitations and what
generalises, and Section~\ref{sec:conclusion} concludes.
